% ---- pg-027 （书页 15） ----

\mvsub{多种二维散布图的联系}

在许多令人兴奋的新的图解法中,有一种是用计算机来联系众多二维散布图.

\mvbold{例 1.8}\quad（散布图的联系及刷亮）

参看\hyperref[tab:1.2]{表 1.2} 的纸张性能数据,我们来说明二维散布图的联系.这些数据表示了三个变量的测量值: $x_1=$ 密度, $x_2=$ 沿机器方向的强度; $x_3=$ 与机器交叉方向的强度.\hyperref[fig:1.9]{图 1.9} 显示了一个 $3\times3$ 阵列,由每对变量做成的二维散布图组成.例如,图中左上角那个图是由 $(x_1, x_2)$ 这一对观测值做成的散布图.即, $x_1$ 的值是沿着水平坐标轴绘图,而 $x_2$ 的值则沿着垂直坐标轴绘图.图中右下角的散布图是由 $(x_2, x_3)$ 这对观测值做成的,将坐标轴颠倒一下就成.图中其他散布图的相应说明也成立.注意变量及它们的二维排列在框图中表示时是沿着西南-东北对角线方向的.将\hyperref[fig:1.9]{图 1.9} 的 $(x_1, x_2)$ 散布图中明显的离群值做出标记(选择出),产生了\hyperref[fig:1.10]{图 1.10}(a),离群值被标为样本 25,同一个数据点在所有的散布图中被突出标示.可发现样本 25 也以离群值出现在 $(x_1, x_3)$ 的散布图中,但不是 $(x_2, x_3)$ 的散布图的离群值.删掉这个样本形成改变后的见\hyperref[fig:1.10]{图 1.10}(b)散布图.

\begin{mvfigure}
\mvinc{pg-027}{0.29268,0.39312,0.72902,0.67298}
\phantomsection\label{fig:1.9}\mvcap{图 1.9\quad 表 1.2 纸张性能数据散布图}
\end{mvfigure}

从\hyperref[fig:1.10]{图 1.10} 中,我们注意到,在 $(x_3, x_c)$ 散布图中一些点看起来与其他的点没有联系.用一个虚线矩形框将这些点框起,在其他所有散布图中高亮显示这些选中的点,如\hyperref[fig:1.11]{图 1.11}(a) 所示.进一步检查显示出样本 16$\sim$21,样本 34 及样本 38$\sim$41 实际上是从一卷旧纸中选取的样本,它们是为了使生产的卡纸板有足够的厚度.删除离群值及与旧纸相关的点,调整剩下的观测值的排列,得到\hyperref[fig:1.11]{图 1.11}(b)中的散布图.

高亮显示对应于一个变量散布图中被选中的区域的点,称为刷亮.刷亮可以从一个矩形开始,如\hyperref[fig:1.11]{图 1.11}(a),然后,刷子可以移动给出一个高亮显示点的序列.这个过程可在任何时候停止,以得到一个当前情况的瞬态.\quad$\blacksquare$

像例 1.8 中的那种散布图对数据分析非常有帮助.另一个新的重要的图解法使用了计算机软件,使数据分析者可用种种不同的三维透视图来观察高维数据.这可以动态地,持续地进%
